% --- Archon physics formalization source begin ---
% archon:physics
% archon:covers PhyXMiniProblems/problem_phyx_mini_0500.lean
% archon:source-report reports/phyx_mini/problem_phyx_mini_0500.source.json

\chapter{Physics problem phyx\_mini\_0500}
\label{ch:PhyXMiniProblems_problem_phyx_mini_0500}

\paragraph{Problem source.}
\#\# Physical scenario

An experiment was performed in which neutrons were shot through two slits spaced \textbackslash{}(0.10\textbackslash{}text\{ nm\}\textbackslash{}) apart and detected \textbackslash{}(3.5\textbackslash{}text\{ m\}\textbackslash{}) behind the slits. Figure shows the detector output. Notice the \textbackslash{}(100\textbackslash{} \textbackslash{}mu\textbackslash{}text\{m\}\textbackslash{}) scale on the figure.

\#\# Question

To one significant figure, what was the speed of the neutrons?

\#\# Answer choices

- A: A: \textbackslash{}( 150\textbackslash{}text\{ m/s\} \textbackslash{})

- B: B: \textbackslash{}( 170\textbackslash{}text\{ m/s\} \textbackslash{})

- C: C: \textbackslash{}( 200\textbackslash{}text\{ m/s\} \textbackslash{})

- D: D: \textbackslash{}( 120\textbackslash{}text\{ m/s\} \textbackslash{})

\#\# Figure caption (auxiliary; use the image as primary evidence)

The image is a graph depicting a waveform representing neutron intensity. 

- **X-axis**: Not labeled with specific values but used as the horizontal reference line.
- **Y-axis**: Labeled as "Neutron intensity," indicating the measured parameter.
- **Waveform**: The shape resembles a series of peaks and troughs, gradually increasing in amplitude towards the center and decreasing symmetrically.
- **Scale Bar**: Above the waveform, there is a scale indicator labeled "100 μm," with a two-headed arrow suggesting the physical scale representation on the graph.

The graph likely shows an interference pattern, possibly from an experiment involving neutron beams and slit diffraction.

\#\# Dataset metadata

Quantum Phenomena; Physical Model Grounding Reasoning, Multi-Formula Reasoning

\paragraph{Recorded answer/context.}
C: C: \textbackslash{}( 200\textbackslash{}text\{ m/s\} \textbackslash{})

\paragraph{Figure/image path.}
/root/proposal\_for\_physic/science-mango/phyx\_mini\_run/phyx\_data/test\_image/500.png

\paragraph{Formalization target.}
create a compiling Lean file with sorry bodies at `PhyXMiniProblems/problem\_phyx\_mini\_0500.lean`. The Lean declarations must preserve the physical quantities, dimensions or dimensional roles, figure labels, governing-law hypotheses, and final relation expressed by this problem.
Use Mathlib/Physlib names found through LeanExplore where available. If a physics API is missing, introduce faithful local abstractions rather than scalar placeholder aliases.

\begin{theorem}[Physics formalization target]
\label{thm:physics:phyx_mini_0500:target}
\lean{PhyXMiniProblems.ProblemPhyXMini0500.problem_phyx_mini_0500}
\uses{def:physics:phyx-mini-0500:phyxminiproblems-problemphyxmini0500-lengthinfemtometers, def:physics:phyx-mini-0500:phyxminiproblems-problemphyxmini0500-lengthinmicrometers, def:physics:phyx-mini-0500:phyxminiproblems-problemphyxmini0500-speedinmeterspersecond, def:physics:phyx-mini-0500:phyxminiproblems-problemphyxmini0500-neutrondoubleslitsetup, def:physics:phyx-mini-0500:phyxminiproblems-problemphyxmini0500-matchesneutrondoubleslitscenario, def:physics:phyx-mini-0500:phyxminiproblems-problemphyxmini0500-matchesproblemlengthreadouts, def:physics:phyx-mini-0500:phyxminiproblems-problemphyxmini0500-matchessupplieddetectorfigure, def:physics:phyx-mini-0500:phyxminiproblems-problemphyxmini0500-usesstandardneutronreferencedata, def:physics:phyx-mini-0500:phyxminiproblems-problemphyxmini0500-hasphysicalneutroninterferenceparameters, def:physics:phyx-mini-0500:phyxminiproblems-problemphyxmini0500-satisfiesneutroninterferencelaws, def:physics:phyx-mini-0500:phyxminiproblems-problemphyxmini0500-answerchoice, def:physics:phyx-mini-0500:phyxminiproblems-problemphyxmini0500-roundstoonesignificantfigure, def:physics:phyx-mini-0500:phyxminiproblems-problemphyxmini0500-matchesdisplayedspeed}
The assigned autoformalize agent should translate this physics problem into Lean declarations in the covered file, with theorem and lemma proof bodies written as `by sorry`.
\end{theorem}
\begin{proof}
This is an autoformalization task, not a proof task. Produce statements that can later be proved by the physics prover without weakening the physical model.
\end{proof}

% --- Archon declaration topology begin ---
\section{Declaration topology}
\begin{definition}[Length Quantity]
  \label{def:physics:phyx-mini-0500:phyxminiproblems-problemphyxmini0500-lengthquantity}
  \lean{PhyXMiniProblems.ProblemPhyXMini0500.LengthQuantity}
  A nonnegative, unit-independent physical quantity with length dimension
\end{definition}
\begin{proof}
  This is the declared model definition of Length Quantity.
\end{proof}
\begin{definition}[Mass Quantity]
  \label{def:physics:phyx-mini-0500:phyxminiproblems-problemphyxmini0500-massquantity}
  \lean{PhyXMiniProblems.ProblemPhyXMini0500.MassQuantity}
  A nonnegative, unit-independent physical mass
\end{definition}
\begin{proof}
  This is the declared model definition of Mass Quantity.
\end{proof}
\begin{definition}[action Dimension]
  \label{def:physics:phyx-mini-0500:phyxminiproblems-problemphyxmini0500-actiondimension}
  \lean{PhyXMiniProblems.ProblemPhyXMini0500.actionDimension}
  The dimension of action, mass * length\textasciicircum{}2 / time
\end{definition}
\begin{proof}
  This is the declared model definition of action Dimension.
\end{proof}
\begin{definition}[Action Quantity]
  \label{def:physics:phyx-mini-0500:phyxminiproblems-problemphyxmini0500-actionquantity}
  \lean{PhyXMiniProblems.ProblemPhyXMini0500.ActionQuantity}
  \uses{def:physics:phyx-mini-0500:phyxminiproblems-problemphyxmini0500-actiondimension}
  A nonnegative, unit-independent physical action
\end{definition}
\begin{proof}
  This is the declared model definition of Action Quantity.
\end{proof}
\begin{definition}[length Readout]
  \label{def:physics:phyx-mini-0500:phyxminiproblems-problemphyxmini0500-lengthreadout}
  \lean{PhyXMiniProblems.ProblemPhyXMini0500.lengthReadout}
  \uses{def:physics:phyx-mini-0500:phyxminiproblems-problemphyxmini0500-lengthquantity}
  Read a physical length as a real number in a selected length unit
\end{definition}
\begin{proof}
  This is the declared model definition of length Readout.
\end{proof}
\begin{definition}[length In Meters]
  \label{def:physics:phyx-mini-0500:phyxminiproblems-problemphyxmini0500-lengthinmeters}
  \lean{PhyXMiniProblems.ProblemPhyXMini0500.lengthInMeters}
  \uses{def:physics:phyx-mini-0500:phyxminiproblems-problemphyxmini0500-lengthquantity, def:physics:phyx-mini-0500:phyxminiproblems-problemphyxmini0500-lengthreadout}
  Metre readout of a physical length
\end{definition}
\begin{proof}
  This is the declared model definition of length In Meters.
\end{proof}
\begin{definition}[length In Femtometers]
  \label{def:physics:phyx-mini-0500:phyxminiproblems-problemphyxmini0500-lengthinfemtometers}
  \lean{PhyXMiniProblems.ProblemPhyXMini0500.lengthInFemtometers}
  \uses{def:physics:phyx-mini-0500:phyxminiproblems-problemphyxmini0500-lengthquantity, def:physics:phyx-mini-0500:phyxminiproblems-problemphyxmini0500-lengthreadout}
  Femtometre readout of a physical length
\end{definition}
\begin{proof}
  This is the declared model definition of length In Femtometers.
\end{proof}
\begin{definition}[length In Nanometers]
  \label{def:physics:phyx-mini-0500:phyxminiproblems-problemphyxmini0500-lengthinnanometers}
  \lean{PhyXMiniProblems.ProblemPhyXMini0500.lengthInNanometers}
  \uses{def:physics:phyx-mini-0500:phyxminiproblems-problemphyxmini0500-lengthquantity, def:physics:phyx-mini-0500:phyxminiproblems-problemphyxmini0500-lengthreadout}
  Nanometre readout of a physical length
\end{definition}
\begin{proof}
  This is the declared model definition of length In Nanometers.
\end{proof}
\begin{definition}[length In Micrometers]
  \label{def:physics:phyx-mini-0500:phyxminiproblems-problemphyxmini0500-lengthinmicrometers}
  \lean{PhyXMiniProblems.ProblemPhyXMini0500.lengthInMicrometers}
  \uses{def:physics:phyx-mini-0500:phyxminiproblems-problemphyxmini0500-lengthquantity, def:physics:phyx-mini-0500:phyxminiproblems-problemphyxmini0500-lengthreadout}
  Micrometre readout of a physical length
\end{definition}
\begin{proof}
  This is the declared model definition of length In Micrometers.
\end{proof}
\begin{definition}[mass In Kilograms]
  \label{def:physics:phyx-mini-0500:phyxminiproblems-problemphyxmini0500-massinkilograms}
  \lean{PhyXMiniProblems.ProblemPhyXMini0500.massInKilograms}
  \uses{def:physics:phyx-mini-0500:phyxminiproblems-problemphyxmini0500-massquantity}
  Kilogram readout of a physical mass
\end{definition}
\begin{proof}
  This is the declared model definition of mass In Kilograms.
\end{proof}
\begin{definition}[speed In Meters Per Second]
  \label{def:physics:phyx-mini-0500:phyxminiproblems-problemphyxmini0500-speedinmeterspersecond}
  \lean{PhyXMiniProblems.ProblemPhyXMini0500.speedInMetersPerSecond}
  Metre-per-second readout of a physical speed
\end{definition}
\begin{proof}
  This is the declared model definition of speed In Meters Per Second.
\end{proof}
\begin{definition}[action In Joule Seconds]
  \label{def:physics:phyx-mini-0500:phyxminiproblems-problemphyxmini0500-actioninjouleseconds}
  \lean{PhyXMiniProblems.ProblemPhyXMini0500.actionInJouleSeconds}
  \uses{def:physics:phyx-mini-0500:phyxminiproblems-problemphyxmini0500-actionquantity}
  SI readout of an action in joule-seconds
\end{definition}
\begin{proof}
  This is the declared model definition of action In Joule Seconds.
\end{proof}
\begin{definition}[Particle Species]
  \label{def:physics:phyx-mini-0500:phyxminiproblems-problemphyxmini0500-particlespecies}
  \lean{PhyXMiniProblems.ProblemPhyXMini0500.ParticleSpecies}
  The particle species sent through the slits
\end{definition}
\begin{proof}
  This is the declared model definition of Particle Species.
\end{proof}
\begin{definition}[Slit Arrangement]
  \label{def:physics:phyx-mini-0500:phyxminiproblems-problemphyxmini0500-slitarrangement}
  \lean{PhyXMiniProblems.ProblemPhyXMini0500.SlitArrangement}
  The aperture arrangement used by the experiment
\end{definition}
\begin{proof}
  This is the declared model definition of Slit Arrangement.
\end{proof}
\begin{definition}[Propagation Regime]
  \label{def:physics:phyx-mini-0500:phyxminiproblems-problemphyxmini0500-propagationregime}
  \lean{PhyXMiniProblems.ProblemPhyXMini0500.PropagationRegime}
  Approximation regime used for the textbook calculation
\end{definition}
\begin{proof}
  This is the declared model definition of Propagation Regime.
\end{proof}
\begin{definition}[Plot Axis]
  \label{def:physics:phyx-mini-0500:phyxminiproblems-problemphyxmini0500-plotaxis}
  \lean{PhyXMiniProblems.ProblemPhyXMini0500.PlotAxis}
  The two axes of the detector-output graph
\end{definition}
\begin{proof}
  This is the declared model definition of Plot Axis.
\end{proof}
\begin{definition}[Plot Axis Role]
  \label{def:physics:phyx-mini-0500:phyxminiproblems-problemphyxmini0500-plotaxisrole}
  \lean{PhyXMiniProblems.ProblemPhyXMini0500.PlotAxisRole}
  Physical role of an axis in the detector-output graph
\end{definition}
\begin{proof}
  This is the declared model definition of Plot Axis Role.
\end{proof}
\begin{definition}[Figure Landmark]
  \label{def:physics:phyx-mini-0500:phyxminiproblems-problemphyxmini0500-figurelandmark}
  \lean{PhyXMiniProblems.ProblemPhyXMini0500.FigureLandmark}
  Horizontal landmarks measured in the supplied raster
\end{definition}
\begin{proof}
  This is the declared model definition of Figure Landmark.
\end{proof}
\begin{definition}[Neutron Detector Figure]
  \label{def:physics:phyx-mini-0500:phyxminiproblems-problemphyxmini0500-neutrondetectorfigure}
  \lean{PhyXMiniProblems.ProblemPhyXMini0500.NeutronDetectorFigure}
  \uses{def:physics:phyx-mini-0500:phyxminiproblems-problemphyxmini0500-lengthquantity, def:physics:phyx-mini-0500:phyxminiproblems-problemphyxmini0500-plotaxis, def:physics:phyx-mini-0500:phyxminiproblems-problemphyxmini0500-plotaxisrole, def:physics:phyx-mini-0500:phyxminiproblems-problemphyxmini0500-figurelandmark}
  Qualitative graph information and rounded horizontal pixel coordinates from the primary image. The scaleBarLength is a physical detector length; the pixel coordinates are dimensionless raster readouts
\end{definition}
\begin{proof}
  This is the declared model definition of Neutron Detector Figure.
\end{proof}
\begin{definition}[Neutron Double Slit Setup]
  \label{def:physics:phyx-mini-0500:phyxminiproblems-problemphyxmini0500-neutrondoubleslitsetup}
  \lean{PhyXMiniProblems.ProblemPhyXMini0500.NeutronDoubleSlitSetup}
  \uses{def:physics:phyx-mini-0500:phyxminiproblems-problemphyxmini0500-lengthquantity, def:physics:phyx-mini-0500:phyxminiproblems-problemphyxmini0500-massquantity, def:physics:phyx-mini-0500:phyxminiproblems-problemphyxmini0500-actionquantity, def:physics:phyx-mini-0500:phyxminiproblems-problemphyxmini0500-particlespecies, def:physics:phyx-mini-0500:phyxminiproblems-problemphyxmini0500-slitarrangement, def:physics:phyx-mini-0500:phyxminiproblems-problemphyxmini0500-propagationregime, def:physics:phyx-mini-0500:phyxminiproblems-problemphyxmini0500-neutrondetectorfigure}
  Independent physical quantities of the experiment. The inferred adjacent fringe spacing, matter wavelength, and neutron speed are fields rather than definitions; the calibration, interference, and de Broglie hypotheses below relate them
\end{definition}
\begin{proof}
  This is the declared model definition of Neutron Double Slit Setup.
\end{proof}
\begin{definition}[Matches Neutron Double Slit Scenario]
  \label{def:physics:phyx-mini-0500:phyxminiproblems-problemphyxmini0500-matchesneutrondoubleslitscenario}
  \lean{PhyXMiniProblems.ProblemPhyXMini0500.MatchesNeutronDoubleSlitScenario}
  \uses{def:physics:phyx-mini-0500:phyxminiproblems-problemphyxmini0500-neutrondoubleslitsetup}
  Qualitative physical scenario and approximation selected by the problem
\end{definition}
\begin{proof}
  This is the declared model definition of Matches Neutron Double Slit Scenario.
\end{proof}
\begin{definition}[Matches Problem Length Readouts]
  \label{def:physics:phyx-mini-0500:phyxminiproblems-problemphyxmini0500-matchesproblemlengthreadouts}
  \lean{PhyXMiniProblems.ProblemPhyXMini0500.MatchesProblemLengthReadouts}
  \uses{def:physics:phyx-mini-0500:phyxminiproblems-problemphyxmini0500-lengthinmeters, def:physics:phyx-mini-0500:phyxminiproblems-problemphyxmini0500-lengthinnanometers, def:physics:phyx-mini-0500:phyxminiproblems-problemphyxmini0500-neutrondoubleslitsetup}
  The two calibrated lengths stated in the prose. This premise contains no fringe spacing, wavelength, neutron speed, or answer choice
\end{definition}
\begin{proof}
  This is the declared model definition of Matches Problem Length Readouts.
\end{proof}
\begin{definition}[Matches Supplied Detector Figure]
  \label{def:physics:phyx-mini-0500:phyxminiproblems-problemphyxmini0500-matchessupplieddetectorfigure}
  \lean{PhyXMiniProblems.ProblemPhyXMini0500.MatchesSuppliedDetectorFigure}
  \uses{def:physics:phyx-mini-0500:phyxminiproblems-problemphyxmini0500-lengthinmicrometers, def:physics:phyx-mini-0500:phyxminiproblems-problemphyxmini0500-neutrondetectorfigure}
  Rounded landmark coordinates measured from the 504 × 280 source raster and the physical label on its scale bar. The maxima used for calibration are adjacent peaks near the center. These are figure/data readouts, not a speed or wavelength conclusion
\end{definition}
\begin{proof}
  This is the declared model definition of Matches Supplied Detector Figure.
\end{proof}
\begin{definition}[Uses Standard Neutron Reference Data]
  \label{def:physics:phyx-mini-0500:phyxminiproblems-problemphyxmini0500-usesstandardneutronreferencedata}
  \lean{PhyXMiniProblems.ProblemPhyXMini0500.UsesStandardNeutronReferenceData}
  \uses{def:physics:phyx-mini-0500:phyxminiproblems-problemphyxmini0500-massinkilograms, def:physics:phyx-mini-0500:phyxminiproblems-problemphyxmini0500-actioninjouleseconds, def:physics:phyx-mini-0500:phyxminiproblems-problemphyxmini0500-neutrondoubleslitsetup}
  Standard reference data used by the calculation. Physlib supplies the reduced Planck constant Constants.ℏ in joule-seconds, so the ordinary Planck action is 2πℏ. The neutron mass uses its standard SI value. Neither datum contains the requested speed
\end{definition}
\begin{proof}
  This is the declared model definition of Uses Standard Neutron Reference Data.
\end{proof}
\begin{definition}[Has Physical Neutron Interference Parameters]
  \label{def:physics:phyx-mini-0500:phyxminiproblems-problemphyxmini0500-hasphysicalneutroninterferenceparameters}
  \lean{PhyXMiniProblems.ProblemPhyXMini0500.HasPhysicalNeutronInterferenceParameters}
  \uses{def:physics:phyx-mini-0500:phyxminiproblems-problemphyxmini0500-lengthinmeters, def:physics:phyx-mini-0500:phyxminiproblems-problemphyxmini0500-massinkilograms, def:physics:phyx-mini-0500:phyxminiproblems-problemphyxmini0500-speedinmeterspersecond, def:physics:phyx-mini-0500:phyxminiproblems-problemphyxmini0500-actioninjouleseconds, def:physics:phyx-mini-0500:phyxminiproblems-problemphyxmini0500-neutrondoubleslitsetup}
  Positivity conditions selecting a physically meaningful experiment
\end{definition}
\begin{proof}
  This is the declared model definition of Has Physical Neutron Interference Parameters.
\end{proof}
\begin{definition}[Satisfies Neutron Interference Laws]
  \label{def:physics:phyx-mini-0500:phyxminiproblems-problemphyxmini0500-satisfiesneutroninterferencelaws}
  \lean{PhyXMiniProblems.ProblemPhyXMini0500.SatisfiesNeutronInterferenceLaws}
  \uses{def:physics:phyx-mini-0500:phyxminiproblems-problemphyxmini0500-lengthreadout, def:physics:phyx-mini-0500:phyxminiproblems-problemphyxmini0500-lengthinmeters, def:physics:phyx-mini-0500:phyxminiproblems-problemphyxmini0500-massinkilograms, def:physics:phyx-mini-0500:phyxminiproblems-problemphyxmini0500-speedinmeterspersecond, def:physics:phyx-mini-0500:phyxminiproblems-problemphyxmini0500-actioninjouleseconds, def:physics:phyx-mini-0500:phyxminiproblems-problemphyxmini0500-neutrondoubleslitsetup}
  The three governing relations used in the calculation: a linear image calibration converts the central peak spacing from pixels into a detector length using the 100 μm bar; adjacent paraxial double-slit maxima obey Δy d = L λ; nonrelativistic neutrons obey de Broglie's relation λ m v = h. The first two relations are stated in arbitrary length units. These laws do not specialize the speed, wavelength, fringe spacing, or answer label to the current problem's result
\end{definition}
\begin{proof}
  This is the declared model definition of Satisfies Neutron Interference Laws.
\end{proof}
\begin{definition}[Answer Choice]
  \label{def:physics:phyx-mini-0500:phyxminiproblems-problemphyxmini0500-answerchoice}
  \lean{PhyXMiniProblems.ProblemPhyXMini0500.AnswerChoice}
  Labels of the four speed choices supplied with the problem
\end{definition}
\begin{proof}
  This is the declared model definition of Answer Choice.
\end{proof}
\begin{definition}[displayed Speed In Meters Per Second]
  \label{def:physics:phyx-mini-0500:phyxminiproblems-problemphyxmini0500-displayedspeedinmeterspersecond}
  \lean{PhyXMiniProblems.ProblemPhyXMini0500.displayedSpeedInMetersPerSecond}
  \uses{def:physics:phyx-mini-0500:phyxminiproblems-problemphyxmini0500-answerchoice}
  Displayed speeds in metres per second
\end{definition}
\begin{proof}
  This is the declared model definition of displayed Speed In Meters Per Second.
\end{proof}
\begin{definition}[recorded Dataset Answer]
  \label{def:physics:phyx-mini-0500:phyxminiproblems-problemphyxmini0500-recordeddatasetanswer}
  \lean{PhyXMiniProblems.ProblemPhyXMini0500.recordedDatasetAnswer}
  \uses{def:physics:phyx-mini-0500:phyxminiproblems-problemphyxmini0500-answerchoice}
  The answer label recorded by the source dataset
\end{definition}
\begin{proof}
  This is the declared model definition of recorded Dataset Answer.
\end{proof}
\begin{definition}[Rounds To One Significant Figure]
  \label{def:physics:phyx-mini-0500:phyxminiproblems-problemphyxmini0500-roundstoonesignificantfigure}
  \lean{PhyXMiniProblems.ProblemPhyXMini0500.RoundsToOneSignificantFigure}
  reported is value rounded to one significant figure when it is a single nonzero decimal digit times a power of ten and ordinary nearest-place rounding of value gives that number
\end{definition}
\begin{proof}
  This is the declared model definition of Rounds To One Significant Figure.
\end{proof}
\begin{definition}[Matches Displayed Speed]
  \label{def:physics:phyx-mini-0500:phyxminiproblems-problemphyxmini0500-matchesdisplayedspeed}
  \lean{PhyXMiniProblems.ProblemPhyXMini0500.MatchesDisplayedSpeed}
  \uses{def:physics:phyx-mini-0500:phyxminiproblems-problemphyxmini0500-speedinmeterspersecond, def:physics:phyx-mini-0500:phyxminiproblems-problemphyxmini0500-neutrondoubleslitsetup, def:physics:phyx-mini-0500:phyxminiproblems-problemphyxmini0500-answerchoice, def:physics:phyx-mini-0500:phyxminiproblems-problemphyxmini0500-displayedspeedinmeterspersecond, def:physics:phyx-mini-0500:phyxminiproblems-problemphyxmini0500-roundstoonesignificantfigure}
  A displayed choice agrees with the modeled speed to one significant figure
\end{definition}
\begin{proof}
  This is the declared model definition of Matches Displayed Speed.
\end{proof}
\begin{lemma}[adjacent fringe spacing micrometers eq]
  \label{lem:physics:phyx-mini-0500:phyxminiproblems-problemphyxmini0500-adjacent-fringe-spacing-micrometers-eq}
  \lean{PhyXMiniProblems.ProblemPhyXMini0500.adjacent_fringe_spacing_micrometers_eq}
  \uses{def:physics:phyx-mini-0500:phyxminiproblems-problemphyxmini0500-lengthinmicrometers, def:physics:phyx-mini-0500:phyxminiproblems-problemphyxmini0500-neutrondoubleslitsetup, def:physics:phyx-mini-0500:phyxminiproblems-problemphyxmini0500-matchessupplieddetectorfigure, def:physics:phyx-mini-0500:phyxminiproblems-problemphyxmini0500-satisfiesneutroninterferencelaws}
  The primary-image calibration first determines an adjacent central-fringe spacing of 100 μm × 32/48 = 200/3 μm
\end{lemma}
\begin{proof}
  The conclusion follows from the governing relations and figure readouts named in the statement of adjacent fringe spacing micrometers eq.
\end{proof}
\begin{lemma}[matter wavelength femtometers eq]
  \label{lem:physics:phyx-mini-0500:phyxminiproblems-problemphyxmini0500-matter-wavelength-femtometers-eq}
  \lean{PhyXMiniProblems.ProblemPhyXMini0500.matter_wavelength_femtometers_eq}
  \uses{def:physics:phyx-mini-0500:phyxminiproblems-problemphyxmini0500-lengthinfemtometers, def:physics:phyx-mini-0500:phyxminiproblems-problemphyxmini0500-neutrondoubleslitsetup, def:physics:phyx-mini-0500:phyxminiproblems-problemphyxmini0500-matchesproblemlengthreadouts, def:physics:phyx-mini-0500:phyxminiproblems-problemphyxmini0500-matchessupplieddetectorfigure, def:physics:phyx-mini-0500:phyxminiproblems-problemphyxmini0500-satisfiesneutroninterferencelaws}
  Combining that calibrated fringe spacing with d = 0.10 nm, L = 3.5 m, and the paraxial double-slit law gives λ = 40/21 fm
\end{lemma}
\begin{proof}
  The conclusion follows from the governing relations and figure readouts named in the statement of matter wavelength femtometers eq.
\end{proof}
\begin{lemma}[neutron speed in literal data rounding interval]
  \label{lem:physics:phyx-mini-0500:phyxminiproblems-problemphyxmini0500-neutron-speed-in-literal-data-rounding-interval}
  \lean{PhyXMiniProblems.ProblemPhyXMini0500.neutron_speed_in_literal_data_rounding_interval}
  \uses{def:physics:phyx-mini-0500:phyxminiproblems-problemphyxmini0500-speedinmeterspersecond, def:physics:phyx-mini-0500:phyxminiproblems-problemphyxmini0500-neutrondoubleslitsetup, def:physics:phyx-mini-0500:phyxminiproblems-problemphyxmini0500-matchesproblemlengthreadouts, def:physics:phyx-mini-0500:phyxminiproblems-problemphyxmini0500-matchessupplieddetectorfigure, def:physics:phyx-mini-0500:phyxminiproblems-problemphyxmini0500-usesstandardneutronreferencedata, def:physics:phyx-mini-0500:phyxminiproblems-problemphyxmini0500-hasphysicalneutroninterferenceparameters, def:physics:phyx-mini-0500:phyxminiproblems-problemphyxmini0500-satisfiesneutroninterferencelaws}
  The literal 0.10 nm slit separation, calibrated pattern, double-slit relation, de Broglie relation, and standard constants put the neutron speed in the nearest-10\textasciicircum{}8 m/s interval that rounds to 2 × 10\textasciicircum{}8 m/s. This is not the nearest-hundred interval associated with the source's recorded choice C
\end{lemma}
\begin{proof}
  The conclusion follows from the governing relations and figure readouts named in the statement of neutron speed in literal data rounding interval.
\end{proof}
% --- Archon declaration topology end ---
% --- Archon physics formalization source end ---
